\documentclass[10pt,a4paper]{amsart}
\usepackage[margin=19mm]{geometry}
\usepackage{amsmath,amssymb,mathtools}
\usepackage[scr=rsfs,cal=euler]{mathalfa}
\usepackage{xfrac}
\usepackage[unicode,hidelinks,bookmarks=false]{hyperref}
\newcommand{\ie}{i.e.\ }
\newcommand{\cf}{cf.\ }
\newcommand{\resp}{respectively\ }
\newcommand{\Subsection}{\subsection}
\providecommand{\nobreakdash}{\nobreak\hbox{-}}
\setlength{\emergencystretch}{2em}
\multlinegap0pt
\usepackage{amssymb}
\usepackage{mathtools}
\usepackage{tikz}
\usepackage{enumerate}
\usepackage{float}
\usepackage{xcolor}
\newtheorem{thm}{Theorem}
\newtheorem{lem}{Lemma}
\newcommand{\BB}{\mathcal B}
\newcommand{\F}{\mathbb{F}}
\newcommand{\colim}{\operatorname*{colim}}
\newcommand{\pa}{\partial}
\newcommand{\wt}{\widetilde}
\title{\bfseries Boundary-weighted barycenter spaces}
\author{Mohameden Ahmedou, Sadok Kallel}
\date{Source: arXiv:2608.22040v1 (CC BY 4.0)}
\begin{document}
\maketitle

\noindent\emph{Source and licence.} \bfseries Boundary-weighted barycenter spaces. Version arXiv:2608.22040v1 (CC BY 4.0), distributed by arXiv under the Creative Commons Attribution 4.0 International licence, http://creativecommons.org/licenses/by/4.0/, as recorded in the arXiv licence metadata for this exact version and shown on the abstract page https://arxiv.org/abs/2608.22040v1. Full licence notice: \texttt{licenses/arxiv-2608.22040-CC-BY-4.0.txt}.\par\smallskip

\noindent\textit{Editorial scope.} Counted unit: the article's thm thm:intro-general-homology (source0.tex lines 155--187). The article's complete original proof of it is delivered with it (source0.tex lines 1466--1575, one \texttt{proof} environment of 110 lines, which the article places away from the statement). No mathematics is written, repaired or re-derived: the statement and the proof are the article's own lines, in the article's own notation, and no retained line is substituted or re-flowed.  Retained uncounted as the source-local context the proof needs: lines 514--524; lines 1288--1290; lines 1364--1373; lines 1378--1389; lines 1405--1418.  Nothing was omitted from any retained span, so the package carries no omission record.  No retained line contains a citation, so the wrapper carries no bibliography and no bibliography entry.  No retained line contains a figure environment, an image-inclusion command or drawing code, so no image is delivered and the package declares no asset dependency.  Editorial formatting only, in addition: an AMS \texttt{amsart} preamble that restates the article's own declarations and macros used by the retained lines, namely the environments \texttt{thm}, \texttt{lem} on separate counters, together with the standard packages those lines need.  The retained environments print in this document as Theorem 1, Lemma 1 in order of appearance, whereas the article prints the same items with the numbering of its own full document.  The complete native source of the article as distributed in the arXiv e-print for this exact version is bundled unchanged as \texttt{sources/208297/source0.tex}.
\begin{lem}\label{lem:compatible-triangulations}
For a fixed order \(l\), the finite family consisting of the spaces
\(\BB_n(M)\), the closed strata \(\BB_q^p(M)\), and their intersections
admits compatible CW structures in which every inclusion occurring in
\(\mathcal D_l\) is the inclusion of a subcomplex. Consequently all these
inclusions are closed cofibrations, and the diagram \(\mathcal D_l\) is Reedy
cofibrant\footnote{Let $P$ be a poset. A diagram $
X\colon P\longrightarrow \mathbf{Top}
$
is Reedy cofibrant if every ``latching map'' $\colim_{\beta<\alpha}X_\beta\rightarrow X_\alpha$ is a  cofibration.}.
\end{lem}

\begin{equation}\label{convention}
X*S^{-1}=S^{-1}*X=X.
\end{equation}

\begin{equation}\label{eq:Ti-quotient-general} \frac{T_l^{(i)}}{T_l^{(i-1)}}\simeq\left(\frac{\BB_{l+i}(\pa M)}{\BB_{l+i-1}(\pa M)}\right)
\vee\bigvee_{1\le j<i}
\left(
\frac{\BB_{l+i-2j}^{j}(M)}
{\BB_{l+i-2j-1}^{j}(M)\cup\BB_{l+i-2j+1}^{j-1}(M)}
\right)
\vee\left(
\frac{\BB_{l-i}^{i}(M)}{\BB_{l-i+1}^{i-1}(M)}
\right)
\end{equation}

\begin{lem}\label{lem:row-quotient-splitting}
For $p\geq 1$, we have the homotopy equivalence 
$\displaystyle\frac{\BB_q^p(M)}{\BB_{q+1}^{p-1}(M)} 
\simeq 
A_p*\BB_q(\pa M)$.
In particular
$$\displaystyle 
\wt H_*\left(\BB_q^p(M), \BB_{q+1}^{p-1}(M)\right)
\cong
\wt H_*(A_p*\BB_q(\pa M)).
$$
\end{lem} 

\begin{lem}\label{eq:two-sided-mixed-quotient}
For \(p\ge1\) and \(q\ge2\), there is a homotopy equivalence
$$\frac{\BB_q^p(M)}{\BB_{q-1}^{p}(M)\cup\BB_{q+1}^{p-1}(M)}\simeq
A_p*(\BB_q(\pa M)\vee\Sigma \BB_{q-1}(\pa M))
$$
In particular
$$\displaystyle 
\wt H_*\left(
\BB_q^p(M), \BB_{q-1}^{p}(M)\cup\BB_{q+1}^{p-1}(M)
\right) \cong
\wt H_*(A_p*\BB_q(\pa M))
\oplus
\wt H_*(A_p*\Sigma \BB_{q-1}(\pa M)).$$
\end{lem}

\begin{thm}\label{thm:intro-general-homology}
Let \(M\) be a compact connected smooth manifold of dimension at least two with nonempty boundary.  For every \(l\ge1\), with field coefficients,
\begin{equation}\label{eq:intro-general-even}
\begin{aligned}
\widetilde H_k\bigl(\BB_{2l}^{\partial}(M);\mathbb F\bigr)
\cong {}&
\widetilde H_k\bigl(\BB_{2l}(\partial M);\mathbb F\bigr)
\oplus
\widetilde H_k\bigl(\BB_l(M/\partial M);\mathbb F\bigr)
\\
&\oplus
\bigoplus_{j=1}^{l-1}
\bigoplus_{\substack{a,b\ge0\\a+b=k-1}}
\widetilde H_a\bigl(\BB_j(M/\partial M);\mathbb F\bigr)
\otimes_{\mathbb F}
\widetilde H_b\bigl(\BB_{2l-2j}(\partial M);\mathbb F\bigr).
\end{aligned}
\end{equation}
\begin{equation}\label{eq:intro-general-odd}
\begin{aligned}
\widetilde H_k\bigl(\BB_{2l-1}^{\partial}(M);\mathbb F\bigr)
\cong {}&
\widetilde H_k\bigl(\BB_{2l-1}(\partial M);\mathbb F\bigr)
\\
&\oplus
\bigoplus_{i=1}^{l-1}
\bigoplus_{\substack{a,b\ge0\\a+b=k-1}}
\widetilde H_a\bigl(\BB_i(M/\partial M);\mathbb F\bigr)
\otimes_{\mathbb F}
\widetilde H_b\bigl(\BB_{2l-2i-1}(\partial M);\mathbb F\bigr).
\end{aligned}
\end{equation}
\end{thm}

\begin{proof}[Proof of Theorem \ref{thm:intro-general-homology}]
We prove the claimed decomposition by induction on \(i\), for
\(0\le i\le l\).
We first rewrite \eqref{eq:Ti-quotient-general} using the homotopy equivalence in Lemma \ref{lem:row-quotient-splitting} and the decomposition in Lemma \ref{eq:two-sided-mixed-quotient}, using distributivity of the join over wedges. This gives
\begin{equation}\label{eq:Ti-quotient-expanded}
\begin{aligned}
\frac{T_l^{(i)}}
     {T_l^{(i-1)}}
\simeq {}&
\BB_{l+i}(\pa M)
\vee
\Sigma\BB_{l+i-1}(\pa M)
\\
&\vee
\bigvee_{1\leq j<i}
\left(
A_j*\BB_{l+i-2j}(\pa M)
\right)
\\
&\vee
\bigvee_{1\leq j<i}
\left(
A_j*\Sigma\BB_{l+i-2j-1}(\pa M)
\right)
\\
&\vee
\left(
A_i*\BB_{l-i}(\pa M)
\right),
\end{aligned}
\end{equation}
The null-homotopy needed at this point preserves the rank filtration.
Indeed, choose \(y_0\in\pa M\) and put
\[
H(s,\sigma)=s\delta_{y_0}+(1-s)\sigma .
\]
If \(\sigma\in\BB_{l+i-1-2j}^{j}(M)\subset T_l^{(i-1)}\), then adding
one boundary slot gives
\(H(s,\sigma)\in\BB_{l+i-2j}^{j}(M)\subset T_l^{(i)}\).
Thus \(T_l^{(i-1)}\) is null-homotopic inside \(T_l^{(i)}\). Since the
inclusion is a cofibration by Lemma \ref{lem:compatible-triangulations},
\[
\frac{T_l^{(i)}}{T_l^{(i-1)}}\simeq
T_l^{(i)}\vee\Sigma T_l^{(i-1)}.
\]
We can therefore rewrite the preceding equivalence with one suspended term
on each side:
\begin{equation}\label{eq:Ti-quotient-regrouped}
\begin{aligned}
 T_l^{(i)}\vee\Sigma T_l^{(i-1)}\simeq {}&
\left[
\BB_{l+i}(\pa M)
\vee
\bigvee_{1\leq j\leq i}
\left(
A_j*\BB_{l+i-2j}(\pa M)
\right)
\right]
\\
&\vee
\Sigma
\left[
\BB_{l+i-1}(\pa M)
\vee
\bigvee_{1\leq j<i}
\left(
A_j*\BB_{l+i-2j-1}(\pa M)
\right)
\right].
\end{aligned}
\end{equation}
If wedge cancellation were valid, induction would identify $T_l^{(i)}(M)$ with the wedge
\begin{equation}\label{splitformula}\BB_{l+i}(\pa M)
\vee
\bigvee_{1\leq j\leq i}
\left(
A_j*\BB_{l+i-2j}(\pa M)
\right)
.\end{equation}
However, cancellation for the wedge product does not hold in general: there are spaces $A,B,C$ such that $A\vee B\simeq A\vee C$, but $B\not\simeq C$. After passing to reduced homology, all groups are finite-dimensional in each degree, and the inductive isomorphism for $T_l^{(i-1)}$ permits degreewise cancellation of the common suspended summand. Thus induction on $i$ gives
\begin{equation}\label{eq:Ti-homology-final}
\wt H_*(T_l^{(i)})
\cong
\wt H_*(\BB_{l+i}(\pa M))
\oplus
\bigoplus_{j=1}^{i}
\wt H_*(\BB_j(M/\pa M)*\BB_{l+i-2j}(\pa M)).
\end{equation}

Putting \(i=l\) yields the formula
$$\wt H_*(\BB_{2l}^{\pa}(M))
\cong
\wt H_*(\BB_{2l}(\pa M))
\oplus
\wt H_*(\BB_l(M/\pa M))
\notag\ \oplus
\bigoplus_{i=1}^{l-1}
\wt H_*(\BB_i(M/\pa M)*\BB_{2l-2i}(\pa M))
$$
Here we use the convention \(\BB_0(\pa M)=S^{-1}\) and $X*S^{-1} = X$ (see \eqref{convention}), which turns the final even mixed term with $l=i=j$ into \(\BB_l(M/\pa M)\), as stated. 

We can finally use the field-coefficient join identity
\begin{equation}\label{eq:join-homology}
\wt H_*(X*Y;\F)
\cong
\Sigma\bigl(\wt H_*(X;\F)\otimes_{\F}\wt H_*(Y;\F)\bigr),
\end{equation}
to obtain the even formula \eqref{eq:intro-general-even}.
Putting \(i=l-1\) in \eqref{eq:Ti-homology-final} and then using \eqref{eq:join-homology} yields the odd formula \eqref{eq:intro-general-odd}.  
\end{proof}

\end{document}
